% Introduction for the BISON manuscript.
% Author: fangtoast
% Email: fangtoast@foxmail.com
% Copyright (c) 2026 fangtoast. All rights reserved.

\section{Introduction}

Two-dimensional photonic crystals control electromagnetic propagation through
the geometry of a periodic unit cell. Periodicity organizes electromagnetic
states into bands and can create frequency intervals in which propagation is
suppressed~\cite{yablonovitch1987inhibited,john1987localization,
joannopoulos2008photonic}. The resulting dispersion relation---the ordered
eigenfrequencies evaluated along a path in the Brillouin zone---determines not
only whether a band gap exists, but also its location, width, neighboring-band
curvature, and slow-light behavior. Designing a useful unit cell is therefore
usually a \emph{full-band} problem rather than a scalar-property regression
problem.

In the forward problem, a structure is given and a Maxwell eigensolver computes
its bands. Plane-wave solvers such as MIT Photonic Bands (MPB) provide a
well-established numerical route~\cite{johnson2001mpb}. In inverse design, a
desired dispersion is given and the structure must be recovered. Classical
topology optimization can search high-dimensional photonic geometries and has
produced low-loss and fabrication-aware devices
~\cite{jensen2004systematic,jensen2011topology,piggott2015demultiplexer,
piggott2017fabrication}. Its repeated forward and adjoint solves, however,
remain costly when many targets or many candidate designs must be considered.

Learning-based inverse design amortizes part of this cost: a trained forward
network can evaluate or differentiate many candidates much faster than a
solver~\cite{peurifoy2018nanophotonic,tahersima2019splitters}. The difficulty is
that the inverse map is one-to-many and the learned objective is imperfect.
Direct regression can average incompatible geometries, whereas a generator or
an optimizer can exploit surrogate error and return a structure that looks
accurate only to the network. Tandem networks address the first issue by
training an inverse model through a frozen forward network
~\cite{liu2018tandem}; time- and budget-aware inverse benchmarks emphasize the
second~\cite{ren2020benchmarking}. Consequently, a persuasive inverse result
must specify both its candidate budget and its independent physical
verification budget.

Neural operators offer a particularly relevant forward representation. They
learn maps between discretized functions, and Fourier neural operators (FNOs)
parameterize global integral kernels through truncated Fourier modes
~\cite{li2021fno,kovachki2023neuraloperator}. This global frequency-domain
mixing matches the periodic input geometry more directly than a flattened
vector regressor. Existing photonic neural operators have primarily targeted
optical fields or device responses~\cite{gu2022neurolight,zhu2024pace,
augenstein2023emno}. Learning the complete eigenfrequency table of a periodic
unit cell, then testing that operator inside a solver-verified inverse loop,
remains a distinct task.

This work introduces \method{} (\textbf{B}and-\textbf{I}nformed
\textbf{S}pectral \textbf{O}perator \textbf{N}etwork), the name we give to
the project's implemented \texttt{fno2d\_\allowbreak band\_\allowbreak v2}
model. BISON maps a periodic permittivity grid directly to the complete sampled
band matrix. Each operator block combines two learned truncated Fourier
kernels, a pointwise path, a data-dependent channel gate, a low-rank global
summary, and a residual update. The full-band output makes BISON useful not
only as a forward predictor, but also as the differentiable objective and
candidate ranker in an inverse-design loop.

Our inverse pipeline deliberately does not use experts. A single-mode latent
conditional decoder samples multiple structures for each target; frozen BISON
scores the candidate pool; an independently trained wave-vector query operator
audits surrogate consistency; and only a fixed Top-5 subset is replayed with
MPB. This separation is central: surrogate accuracy supports efficient search,
whereas real-MPB replay is the only evidence used for physical inverse-design
claims.

Our contributions are fourfold:
\begin{itemize}
  \item We develop BISON, a compact FNO-derived architecture specialized for
  complete photonic band matrices. Its two-scale spectral mixing, adaptive
  gate, low-rank global branch, and residual path are documented as one
  reproducible BandBlock rather than presented under a generic FNO label.
  \item We integrate the frozen operator into an auditable, no-expert inverse
  workflow with equal candidate and solver budgets, a separate query-model
  diagnostic, and per-target solver replay.
  \item We organize evaluation around three questions: whether BISON improves
  held-out full-band prediction, whether coordinate-conditioned query models
  interpolate and extrapolate reliably, and whether BISON-ranked inverse
  candidates survive an identical MPB replay budget.
  \item We report evidence at its actual strength: a formal three-seed forward
  comparison on 10,000 structures, query interpolation and extrapolation
  tests, and a seed-0 inverse replay over the complete 625-target test split.
  BISON reaches \(1.0299\%\pm0.0545\%\) forward MRE; the inverse loop reaches
  \(2.2614\%\) real-MPB MRE, while cross-seed inverse stability remains open.
\end{itemize}

% BISON architecture and solver-verified inverse-design workflow.
% Author: fangtoast
% Email: fangtoast@foxmail.com
% Copyright (c) 2026 fangtoast. All rights reserved.

\begin{figure*}[t]
  \centering
  \resizebox{\textwidth}{!}{%
  \begin{tikzpicture}[
    node distance=5mm and 7mm,
    box/.style={draw, rounded corners=2pt, align=center, minimum height=10mm,
      minimum width=24mm, inner sep=4pt, font=\small},
    data/.style={box, fill=blue!9, draw=blue!55!black},
    operator/.style={box, fill=violet!10, draw=violet!60!black},
    inverse/.style={box, fill=orange!11, draw=orange!65!black},
    solver/.style={box, fill=green!10, draw=green!55!black},
    note/.style={align=center, font=\scriptsize, text=black!75},
    arrow/.style={-{Latex[length=2.2mm]}, thick},
    audit/.style={-{Latex[length=2.2mm]}, thick, dashed, draw=black!60}
  ]
    \node[data] (eps) {periodic permittivity\\\(\permittivity(\mathbf{x})\)};
    \node[operator, right=of eps] (lift) {\(1\!\times\!1\) lift\\\(1\rightarrow32\) channels};
    \node[operator, right=of lift, minimum width=43mm] (blocks) {four BISON BandBlocks\\dual Fourier kernels $+$ pointwise map\\band gate $+$ low-rank global branch $+$ residual};
    \node[operator, right=of blocks] (pool) {global pooling\\and linear head};
    \node[data, right=of pool] (bands) {full band matrix\\\(\bandpred\in\mathbb{R}^{N_k\times N_b}\)};

    \draw[arrow] (eps) -- (lift);
    \draw[arrow] (lift) -- (blocks);
    \draw[arrow] (blocks) -- (pool);
    \draw[arrow] (pool) -- (bands);
    \node[note, above=2mm of blocks] {forward operator: structure field $\rightarrow$ complete sampled dispersion};

    \node[data, below=17mm of eps] (target) {target bands\\\(\bandtarget\)};
    \node[inverse, right=of target] (gen) {single-mode latent decoder\\\(\generator(\bandtarget,z)\)\\no expert routing};
    \node[inverse, right=of gen] (cand) {\(K=128\) candidates\\legality and duplicate checks};
    \node[operator, right=of cand] (rank) {frozen BISON\\primary objective and ranker};
    \node[operator, right=of rank] (query) {query-operator audit\\physical \(\mathbf{k}\) consistency};
    \node[solver, right=of query] (replay) {Top-5 real-MPB replay\\solver-verified bands};

    \draw[arrow] (target) -- (gen);
    \draw[arrow] (gen) -- (cand);
    \draw[arrow] (cand) -- (rank);
    \draw[arrow] (rank) -- (query);
    \draw[arrow] (query) -- (replay);
    \draw[audit] (bands.south) |- (rank.north);
    \node[note, below=2mm of cand] {surrogate quantities guide search; only replay quantities support physical claims};
  \end{tikzpicture}}
  \caption{BISON and the evidence-separated inverse-design loop. The upper row
  is the implemented full-band FNO-derived surrogate (registered in the code as
  \texttt{fno2d\_band\_v2}). The lower row uses the frozen operator to guide a
  no-expert latent candidate generator. The independent query model is an audit
  rather than a second equal-weight ranker, and MPB remains outside training.}
  \Description{Two-row workflow diagram. The upper row maps a periodic
  permittivity field through a lift, four BISON spectral blocks, pooling, and a
  linear head to a full band matrix. The lower row maps target bands through a
  single-mode latent generator, candidate checks, frozen BISON ranking, a query
  audit, and Top-5 MPB replay.}
  \label{fig:architecture}
\end{figure*}

